\documentclass[12pt]{article}
\usepackage{amsmath, amssymb}
\usepackage[margin=1in]{geometry}
\setlength{\parskip}{6pt}
\title{QUBO Formulation: Alkane Structural Isomer Search Toy Problem}
\date{}
\begin{document}
\maketitle

\subsection*{Alkane Structural Isomer Search Toy Problem}

\paragraph{Problem Statement.}
Given a set of $N$ carbon atoms and a set of $M = N(N-1)/2$ candidate carbon-carbon bonds, the goal is to select a subset of bonds that forms a valid alkane carbon skeleton. Validity is enforced through degree constraints specifying the target valence for each atom and a connectivity constraint ensuring the selected bonds form a single connected component (a tree structure). The task is to find any feasible bond selection satisfying these structural constraints.

\paragraph{Decision Variables.}
For each candidate bond $b \in \mathcal{B}$, where $|\mathcal{B}| = M$, introduce a binary decision variable $x_b \in \{0,1\}$. The variable $x_b = 1$ indicates that bond $b$ is selected in the carbon skeleton, and $x_b = 0$ otherwise.

\paragraph{QUBO Derivation.}
Step 1 --- The problem is a pure feasibility search, so the objective function is constant:
\[ \min_{\mathbf{x} \in \{0,1\}^M} 0. \]

Step 2 --- The constraints are expressed in general algebraic form as:
\begin{align}
\sum_{b \in \mathcal{B}} A_{i,b} x_b &= d_i, \quad \forall i \in \{0, \dots, N-1\}, \\
\sum_{b \in \mathcal{B}} x_b &= N - 1,
\end{align}
where $A_{i,b} \in \{0,1\}$ is the incidence matrix indicating whether bond $b$ is incident to atom $i$, and $d_i$ is the target degree for atom $i$.

Step 3 --- We convert each constraint into a quadratic penalty term with weight $\lambda > 0$.
(a) The squared penalty forms are:
\[ P_{\text{deg}} = \lambda \sum_{i=0}^{N-1} \left( \sum_{b \in \mathcal{B}} A_{i,b} x_b - d_i \right)^2, \quad P_{\text{conn}} = \lambda \left( \sum_{b \in \mathcal{B}} x_b - (N-1) \right)^2. \]
(b) Expanding the degree penalty using the idempotent property $x_b^2 = x_k$ and $A_{i,b}^2 = A_{i,b}$:
\begin{align}
\left( \sum_{b} A_{i,b} x_b - d_i \right)^2 &= \sum_{b} A_{i,b}^2 x_b^2 - 2d_i \sum_{b} A_{i,b} x_b + d_i^2 + \sum_{b, b' \in \mathcal{B}, b \neq b'} A_{i,b} A_{i,b'} x_b x_{b'} \nonumber \\
&= \sum_{b \in \mathcal{B}} A_{i,b}(1 - 2d_i) x_b + d_i^2 + \sum_{b, b' \in \mathcal{B}, b \neq b'} A_{i,b} A_{i,b'} x_b x_{b'}.
\end{align}
Expanding the connectivity penalty similarly:
\[ \left( \sum_{b} x_b - (N-1) \right)^2 = \sum_{b \in \mathcal{B}} (1 - 2(N-1)) x_b + (N-1)^2 + \sum_{b, b' \in \mathcal{B}, b \neq b'} x_b x_{b'}. \]
(c) The resulting general penalty contributions to the Hamiltonian are:
\begin{align}
P_{\text{deg}} &= \lambda \sum_{i=0}^{N-1} \left[ \sum_{b \in \mathcal{B}} A_{i,b}(1 - 2d_i) x_b + d_i^2 + \sum_{b, b' \in \mathcal{B}, b \neq b'} A_{i,b} A_{i,b'} x_b x_{b'} \right], \\
P_{\text{conn}} &= \lambda \left[ \sum_{b \in \mathcal{B}} (1 - 2(N-1)) x_b + (N-1)^2 + \sum_{b, b' \in \mathcal{B}, b \neq b'} x_b x_{b'} \right].
\end{align}

Step 4 --- Combining the objective and all penalty terms into the single QUBO Hamiltonian $H(\mathbf{x})$:
\[ H(\mathbf{x}) = P_{\text{deg}} + P_{\text{conn}}. \]
Collecting linear terms for each $x_b$ and quadratic terms for each pair $b < b'$ (noting that $\sum_{b \neq b'}$ becomes $2\sum_{b < b'}$ in the standard QUBO form):
\[ H(\mathbf{x}) = \sum_{b \in \mathcal{B}} \left( \lambda \sum_{i=0}^{N-1} A_{i,b}(1 - 2d_i) + \lambda(1 - 2(N-1)) \right) x_b + \sum_{b, b' \in \mathcal{B}, b < b'} 2\lambda \left( \sum_{i=0}^{N-1} A_{i,b} A_{i,b'} + 1 \right) x_b x_{b'} + c. \]

Step 5 --- Reading off the QUBO matrix entries and offset $c$ in general closed form:
\begin{align}
Q_{b,b} &= \lambda \left[ \sum_{i=0}^{N-1} A_{i,b}(1 - 2d_i) + (1 - 2(N-1)) \right], \quad \forall b \in \mathcal{B}, \\
Q_{b,b'} &= 2\lambda \left[ \sum_{i=0}^{N-1} A_{i,b} A_{i,b'} + 1 \right], \quad \forall b, b' \in \mathcal{B}, b < b', \\
c &= \lambda \left[ \sum_{i=0}^{N-1} d_i^2 + (N-1)^2 \right].
\end{align}

\paragraph{Penalty Weight Justification.}
The penalty weight $\lambda$ must be chosen sufficiently large to ensure that any constraint violation increases the energy above the minimum feasible energy. Since the objective function is constant, the minimum feasible energy is exactly $c$. Any violation of the degree or connectivity constraints introduces a strictly positive penalty term. Specifically, because the variables are binary and the constraints are linear with integer coefficients, the minimum non-zero violation magnitude is $1$, yielding a penalty of at least $\lambda$. Thus, choosing $\lambda > 0$ guarantees that feasible configurations have strictly lower energy than infeasible ones. In practice, $\lambda$ is scaled relative to the problem parameters to maintain numerical stability while preserving the ground state correspondence.

\paragraph{Correctness Argument.}
Minimizing $H(\mathbf{x})$ over $\mathbf{x} \in \{0,1\}^M$ recovers the optimal solution because the QUBO is constructed such that $H(\mathbf{x}) = c$ if and only if $\mathbf{x}$ satisfies all constraints, and $H(\mathbf{x}) > c$ otherwise. Since the objective function is constant, any feasible assignment yields the same energy $c$. Infeasible assignments incur a positive penalty, strictly increasing their energy. Therefore, the global minimum of $H$ corresponds exactly to the set of feasible bond selections, and optimization algorithms targeting the ground state will return valid alkane skeletons.

\paragraph{Illustrative Example.}
Consider a concrete instance with $N=5$ carbon atoms. The number of candidate bonds is $M = 5(4)/2 = 10$. Let the target degrees be $d_0, d_1, d_2, d_3, d_4$. Setting the penalty weight to $\lambda = 100$, the general formulas instantiate as follows. The linear coefficients for each bond $b$ become:
\[ Q_{b,b} = 100 \left[ \sum_{i \in \text{endpoints}(b)} (1 - 2d_i) - 7 \right]. \]
The quadratic coefficients for each pair of distinct bonds $b < b'$ become:
\[ Q_{b,b'} = 200 \left[ \mathbb{I}(\text{bond } b \text{ and } b' \text{ share a node}) + 1 \right]. \]
The constant offset is:
\[ c = 100 \left( d_0^2 + d_1^2 + d_2^2 + d_3^2 + d_4^2 + 16 \right). \]
For a specific target degree sequence $d = [1, 1, 1, 1, 4]$ (corresponding to a star topology), the offset evaluates to $c = 100(1+1+1+1+16) = 2000$. The ground state energy is $2000$, achieved by the bond set forming the star graph. Any deviation from this topology violates at least one constraint, increasing the energy by at least $100$. This concrete instantiation demonstrates how the general algebraic expressions directly map to the QUBO parameters for a fixed molecular graph size.

\end{document}